这一节后面会反复用到高斯分布，先把密度和采样写法放在一起。
若 $X\in\mathbb R^d$，均值为 $\mu$，协方差 $\Sigma$ 正定，则
\[
p(x)=\frac{\exp[-\tfrac12(x-\mu)^\top\Sigma^{-1}(x-\mu)]}
 {(2\pi)^{d/2}|\Sigma|^{1/2}}.
\]
当 $\Sigma=LL^\top$ 时，可以先抽取 $\epsilon\sim\mathcal N(0,I)$，再令
$X=\mu+L\epsilon$。这就是高斯情形的重参数化：随机性放进与参数无关的
$\epsilon$，均值与尺度通过确定性运算作用在它上面。
若 $\mu,L$ 来自神经网络，满足交换求导与期望的条件时，就可以沿这条运算链传播参数梯度。
后面的 DDPM 用同一个技巧直接构造任意时刻的带噪样本。

上面的 $H(X)$ 是离散熵；连续变量对应微分熵
$h(X)=-\int p(x)\log p(x)\,\mathrm dx$，两者不要直接混用。
比如高斯微分熵为 $\tfrac12\log[(2\pi e)^d|\Sigma|]$，可以取负值。
KL 散度则在相应绝对连续条件下保持非负，且通常不对称。
对于两个协方差相同的各向同性高斯，有
\[
D_{\rm KL}\bigl(\mathcal N(\mu_1,\nu^2I)\|
 \mathcal N(\mu_2,\nu^2I)\bigr)
=\frac{\|\mu_1-\mu_2\|^2}{2\nu^2}.
\]
这也预先说明了，为什么后面高斯后验之间的 KL 能转化为均值的平方误差。
